\chapter{Matriks}\label{ch:b1:matrices}

Sebuah matriks adalah \omterm{def:b1:linmaps:def}{pemetaan linear} yang ditulis dalam \omterm{prop:b1:vspaces:coordinates}{koordinat}. Bab ini
menyiapkan kamusnya --- komposisi menjadi hasil kali matriks, kebijektifan
menjadi keterbalikan, dan \omterm{def:b1:matrices:changeofbasis}{perubahan basis} menjadi konjugasi --- beserta
sisi algoritmisnya: yaitu \omterm{met:b1:matrices:gauss}{operasi baris}, perhitungan rank dan
invers. Setelah pertama kali ditemui pada jilid Sekolah Menengah, matriks kini
berpijak pada teori
\cref{ch:b1:vspaces,ch:b1:findim,ch:b1:linmaps}.

\section{Matriks dan pemetaan linear}

\begin{definition}\label{def:b1:matrices:def}
\omterm{def:b1:logic:sets}{Himpunan} $\mathcal{M}_{n,p}(K)$ adalah \omterm{def:b1:vspaces:def}{ruang vektor} berisi larik $n \times p$ yaitu $A
= (a_{ij})$ berisi skalar ($i$: barisnya, $j$: kolomnya), yang berdimensi $np$
(dengan basisnya: matriks $E_{ij}$ yang bernilai $1$ tunggal). Diberikan \omterm{def:b1:vspaces:free}{basis}
$\mathcal{B} = (e_1, \dots, e_p)$ pada $E$ dan $\mathcal{C}$ pada $F$
(dengan $\dim F = n$), \emph{matriks $u \in \mathcal{L}(E, F)$} adalah
larik yang kolom ke-$j$-nya mendaftar \omterm{prop:b1:vspaces:coordinates}{koordinat} $u(e_j)$ pada
$\mathcal{C}$:
\[
\operatorname{Mat}_{\mathcal{B},\mathcal{C}}(u) = (a_{ij}),
\qquad u(e_j) = \sum_{i=1}^{n} a_{ij}\, f_i .
\]
Adapun \omterm{def:b1:logic:map}{pemetaan} $u \mapsto \operatorname{Mat}_{\mathcal{B},\mathcal{C}}(u)$
merupakan isomorfisma dari $\mathcal{L}(E, F)$ pada
$\mathcal{M}_{n,p}(K)$ (\cref{prop:b1:linmaps:basis}: karena \omterm{def:b1:linmaps:def}{pemetaan linear}
tepat merupakan pemilihan peta $e_j$-nya).
\end{definition}

\begin{example}[Turunan, sebagai sebuah matriks]
\label{ex:b1:matrices:derivative}
Misalkan $D(P) = P'$ pada $\R_3[X]$. Pada \omterm{def:b1:vspaces:free}{basis} monomialnya $(1, X, X^2,
X^3)$: $D(1) = 0$, $D(X) = 1$, $D(X^2) = 2X$, $D(X^3) = 3X^2$, sehingga
\[
\operatorname{Mat}(D) =
\begin{pmatrix}
0 & 1 & 0 & 0\\
0 & 0 & 2 & 0\\
0 & 0 & 0 & 3\\
0 & 0 & 0 & 0
\end{pmatrix}.
\]
Sedangkan pada \omterm{def:b1:vspaces:free}{basis} yang \emph{terbagi} $\bigl(1,\ X,\ \frac{X^2}2,\
\frac{X^3}6\bigr)$, setiap vektor basisnya terpetakan ke yang sebelumnya
(karena $D\bigl(\frac{X^k}{k!}\bigr) = \frac{X^{k-1}}{(k-1)!}$), dan
matriksnya menjadi geseran murni: yaitu satuan pada superdiagonalnya, dan nol
di tempat lain. Ada dua moralnya: bahwa matriksnya milik \emph{pasangan}
(\omterm{def:b1:logic:map}{pemetaan}, \omterm{def:b1:vspaces:free}{basis}), bukan \omterm{def:b1:logic:map}{pemetaannya} belaka; dan bahwa \omterm{def:b1:vspaces:free}{basis} yang baik membuat
strukturnya kasatmata sekali pandang --- karena bentuk geserannya seketika menunjukkan
bahwa $D^4 = 0$ pada $\R_3[X]$, dengan setiap pangkat matriksnya mendorong
diagonal satuannya satu langkah lebih jauh ke luar.
\end{example}

\begin{definition}[Hasil kali]\label{def:b1:matrices:product}
Untuk $A \in \mathcal{M}_{n,p}$ dan $B \in \mathcal{M}_{p,q}$:
\[
(AB)_{ik} = \sum_{j=1}^{p} a_{ij}\, b_{jk}
\qquad (1 \leq i \leq n,\ 1 \leq k \leq q).
\]
Dan inilah persis matriks komposisinya:
$\operatorname{Mat}(v \circ u) = \operatorname{Mat}(v)\,
\operatorname{Mat}(u)$ (dengan basisnya cocok di tengah). Demikian pula,
jika $X$ kolom \omterm{prop:b1:vspaces:coordinates}{koordinat} $x$, maka kolom $u(x)$ adalah
$AX$.
\end{definition}

\begin{proof}[Bukti rumus komposisinya]
\[
v(u(e_k)) = v\Bigl(\sum_j b_{jk} f_j\Bigr) = \sum_j b_{jk}\, v(f_j)
= \sum_j b_{jk} \sum_i a_{ij}\, g_i
= \sum_i \Bigl(\sum_j a_{ij} b_{jk}\Bigr) g_i . \qedhere
\]
\end{proof}

\begin{proposition}[Aljabar $\mathcal{M}_n(K)$]\label{prop:b1:matrices:ring}
Matriks persegi $\mathcal{M}_n(K)$ membentuk \omterm{def:b1:structures:ring}{ring} (yang takkomutatif untuk $n
\geq 2$), dengan identitasnya $I_n$; dan \omterm{def:b1:structures:group}{grup} satuannya adalah
\emph{grup linear umum}\index{grup linear umum}
$GL_n(K)$, yang bersesuaian dengan endomorfisma yang \omterm{def:b1:logic:inj}{bijektif}. Untuk $A, B \in
\mathcal{M}_n(K)$:
\[
AB = I_n \implies A \in GL_n(K) \text{ dan } B = A^{-1}
\]
(karena invers sepihaknya menjadi dua sisi, menurut
\cref{cor:b1:linmaps:samedim}).
\end{proposition}

\begin{proof}
Aksioma \omterm{def:b1:structures:ring}{ringnya} terangkut dari $\mathcal{L}(E)$ lewat
isomorfisma pada \cref{def:b1:matrices:def}: karena ia mengubah
komposisi menjadi hasil kali dan jumlah menjadi jumlah, sehingga keasosiatifan,
kedistributifan dan peran $I_n$ terwarisi dari fakta yang
bersesuaian tentang \omterm{def:b1:logic:map}{pemetaannya}, tanpa pemeriksaan entri demi entri.
Adapun ketakkomutatifannya: $E_{12}E_{21} = E_{11} \neq E_{22}
= E_{21}E_{12}$. Jika $AB = I_n$: maka endomorfisma $a$ milik $A$
memenuhi $a \circ b = \mathrm{id}$, sehingga $a$ bersifat \omterm{def:b1:logic:inj}{surjektif} (karena $x =
a(b(x))$ menunjukkan sebuah \omterm{def:b1:logic:map}{prapeta} bagi setiap $x$), jadi \omterm{def:b1:logic:inj}{bijektif} dalam
\omterm{def:b1:findim:def}{dimensi hingga} (\cref{cor:b1:linmaps:samedim}); lalu menyusun $a
\circ b = \mathrm{id}$ dengan $a^{-1}$ di kiri memberikan $b =
a^{-1}$, dan lalu $b\circ a = \mathrm{id}$ pula: sehingga invers sepihaknya
sejak awal dua sisi --- yaitu bantuan yang tegas
\omterm{def:b1:findim:def}{berdimensi hingga}.
\end{proof}

\begin{definition}[Transpos; trace]\label{def:b1:matrices:transpose}
\emph{Transpos} $A = (a_{ij}) \in \mathcal{M}_{n,p}$ adalah $A^{\mathsf T}
= (a_{ji}) \in \mathcal{M}_{p,n}$; yang memenuhi $(AB)^{\mathsf T} = B^{\mathsf T}
A^{\mathsf T}$ dan $(A^{\mathsf T})^{\mathsf T} = A$. Sedangkan
\emph{trace}\index{trace} sebuah matriks persegi adalah $\operatorname{tr} A
= \sum_i a_{ii}$; yang bersifat \omterm{def:b1:linmaps:def}{linear}, dan
\[
\operatorname{tr}(AB) = \operatorname{tr}(BA)
\qquad (A \in \mathcal{M}_{n,p},\ B \in \mathcal{M}_{p,n}).
\]
\end{definition}

\begin{proof}[Bukti kesamaan tracenya]
$\operatorname{tr}(AB) = \sum_i \sum_j a_{ij} b_{ji}$ dan
$\operatorname{tr}(BA) = \sum_j \sum_i b_{ji} a_{ij}$: yaitu jumlah
rangkap yang sama.
\end{proof}

\begin{example}[Trace dalam kerja]
\label{ex:b1:matrices:tracework}
\omterm{def:b1:linmaps:projection}{Proyeksi} pada \cref{ch:b1:linmaps} ke
$\operatorname{Vect}(1,1)$ sepanjang $\operatorname{Vect}(0,1)$,
yaitu $p(x, y) = (x, x)$, bermatriks $A = \begin{pmatrix} 1 & 0\\ 1 &
0\end{pmatrix}$ pada \omterm{def:b1:vspaces:free}{basis} kanoniknya: memang $A^2 = A$, dan
\[
\operatorname{tr} A = 1 = \operatorname{rk} A ,
\]
yang memperlihatkan \cref{exo:b1:matrices:8}: bahwa bagi yang idempoten \omterm{def:b1:matrices:transpose}{tracenya}
\emph{mencacah} dimensi petanya, pada \omterm{def:b1:vspaces:free}{basis} semiring apa pun
matriksnya ditulis. Adapun mekanisme kekekalannya adalah
kesamaan $\operatorname{tr}(AB) = \operatorname{tr}(BA)$:
\[
\operatorname{tr}\bigl(P^{-1}(AP)\bigr) =
\operatorname{tr}\bigl((AP)P^{-1}\bigr) = \operatorname{tr} A ,
\]
sehingga semua matriks yang serupa dengan $A$ berbagi \omterm{def:b1:matrices:transpose}{tracenya} --- yaitu
\emph{invarian numerik} yang pertama bagi sebuah endomorfisma, yang akan disusul
oleh determinan pada \cref{ch:b1:det} (yakni pasangan $(s, p)$ pada
soal akhir pekan di bawah).
\end{example}

\begin{example}[Setangkup tambah antisetangkup]
\label{ex:b1:matrices:symsplit}
Sebutlah $A$ \emph{setangkup} bila $A^{\mathsf T} = A$, dan
\emph{antisetangkup} bila $A^{\mathsf T} = -A$. Maka setiap matriks
persegi terbelah secara tunggal menjadi yang satu tambah yang lain:
\[
A = \underbrace{\frac{A + A^{\mathsf T}}{2}}_{\text{setangkup}}
+ \underbrace{\frac{A - A^{\mathsf T}}{2}}_{\text{antisetangkup}},
\]
dan matriks yang sekaligus keduanya bernilai nol (karena $A = -A$): sehingga kedua \omterm{def:b1:logic:sets}{himpunannya}
merupakan \omterm{def:b1:vspaces:sum}{subruang yang saling melengkapi} pada $\mathcal{M}_n(K)$ --- yaitu
padanan yang persis bagi pembelahan genap/ganjil fungsi
(\cref{ex:b1:vspaces:evenodd}), dengan transposisi memainkan
peran $x \mapsto -x$. Adapun dimensinya: matriks setangkup bersifat \omterm{def:b1:vspaces:free}{bebas}
pada dan di atas diagonalnya, sedangkan yang antisetangkup tegas di atasnya
(karena diagonalnya nol):
\[
\frac{n(n+1)}{2} + \frac{n(n-1)}{2} = n^2 ,
\]
dan cacah yang berimbang itu merupakan penegasan Grassmann atas
kelangsungannya. Untuk $n = 2$: $\begin{pmatrix} 1 & 5\\ 1 &
2\end{pmatrix} = \begin{pmatrix} 1 & 3\\ 3 & 2\end{pmatrix} +
\begin{pmatrix} 0 & 2\\ -2 & 0\end{pmatrix}$. Adapun matriks
setangkup kembali sebagai data \omterm{def:b1:derivative:def}{turunan} kedua pada
\cref{ch:b1:multivar} (yaitu tripel Monge $r, s, t$), sedangkan
yang setangkup sekaligus ortogonal digolongkan pada
\cref{exo:b1:euclid:12}.
\end{example}

\section{Perubahan basis}

\begin{definition}\label{def:b1:matrices:changeofbasis}
Misalkan $\mathcal{B}, \mathcal{B}'$ \omterm{def:b1:vspaces:free}{basis} pada $E$. Maka \emph{matriks
perubahan basis}\index{perubahan basis} $P =
P_{\mathcal{B}\to\mathcal{B}'}$ berkolomkan \omterm{prop:b1:vspaces:coordinates}{koordinat}
vektor \omterm{def:b1:vspaces:free}{basis} yang \emph{baru} pada \omterm{def:b1:vspaces:free}{basis} yang \emph{lama}. Ia
dapat dibalik, $P^{-1} = P_{\mathcal{B}'\to\mathcal{B}}$, dan
\omterm{prop:b1:vspaces:coordinates}{koordinatnya} berubah lewat $X = PX'$ (yaitu lama $=$ $P\,\cdot$ baru).
\end{definition}

\begin{example}[Membaca matriks perubahan basisnya]
\label{ex:b1:matrices:changerun}
Di dalam $\R^2$, dari $\mathcal B$ yang kanonik ke $\mathcal B' =
\bigl((1,1), (1,-1)\bigr)$:
\[
P = P_{\mathcal B\to\mathcal B'} =
\begin{pmatrix} 1 & 1\\ 1 & -1 \end{pmatrix}
\]
(dengan vektor barunya ditulis dalam \omterm{prop:b1:vspaces:coordinates}{koordinat} lamanya, kolom demi kolom). Maka
vektor berkoordinat lama $X = (3, 1)^{\mathsf T}$ berkoordinat
baru $X' = P^{-1}X = \frac12(3 + 1,\ 3 - 1)^{\mathsf T} =
(2, 1)^{\mathsf T}$: memang $2(1,1) + 1(1,-1) = (3,1)$. Awas arahnya
--- karena matriks $P$ dibangun dari \omterm{def:b1:vspaces:free}{basis} yang \emph{baru}
tetapi mengubah \omterm{prop:b1:vspaces:coordinates}{koordinat} \emph{baru menjadi lama} ($X = PX'$); sedangkan beralih
dari lama ke baru berongkos inversnya. Jadi menuliskan periksa kewarasannya
$2(1,1) + (1,-1) = (3,1)$ setelah setiap pengalihan menangkap
galat $P$ yang terbalik, yang merupakan kesalahan paling lazim pada
bab ini.
\end{example}

\begin{theorem}[Perubahan basis bagi sebuah pemetaan]\label{thm:b1:matrices:conjugation}
Misalkan $u \in \mathcal{L}(E)$ bermatriks $A$ pada $\mathcal{B}$ dan $A'$
pada $\mathcal{B}'$, dan $P = P_{\mathcal{B}\to\mathcal{B}'}$. Maka
\[
A' = P^{-1} A\, P .
\]
Dua matriks yang berkaitan dengan cara ini disebut \emph{serupa}\index{matriks
yang serupa}. (Adapun untuk $u \colon E \to F$ dengan dua pasang \omterm{def:b1:vspaces:free}{basis},
rumusnya adalah $A' = Q^{-1} A P$ --- yaitu matriks yang \emph{setara}
satu sama lain\index{matriks yang setara}.)
\end{theorem}

\begin{proof}
Untuk sebarang $x$: $X = PX'$ dan petanya memenuhi $Y = AX$, $Y = PY'$.
Sehingga $PY' = APX'$, yakni $Y' = (P^{-1}AP)X'$ untuk setiap $X'$: jadi matriks
$u$ pada \omterm{def:b1:vspaces:free}{basis} barunya adalah $P^{-1}AP$ (ambillah kolom kanoniknya
bagi $X'$).
\end{proof}

\begin{example}[Basis yang baik membuat sebuah pemetaan bening]
\label{ex:b1:matrices:conjugationrun}
Misalkan $u(x, y) = (y, x)$ (yaitu penukaran), yang bermatriks $A = \begin{pmatrix} 0
& 1\\ 1 & 0\end{pmatrix}$ pada \omterm{def:b1:vspaces:free}{basis} kanoniknya. Adapun pada \omterm{def:b1:vspaces:free}{basis}
$\mathcal B' = \bigl((1,1), (1,-1)\bigr)$:
\[
P = \begin{pmatrix} 1 & 1\\ 1 & -1 \end{pmatrix},
\qquad
P^{-1} = \frac12\begin{pmatrix} 1 & 1\\ 1 & -1 \end{pmatrix},
\qquad
P^{-1} A P = \begin{pmatrix} 1 & 0\\ 0 & -1 \end{pmatrix}.
\]
Sesungguhnya tak ada hasil kali matriks yang diperlukan: karena $u$ menetapkan $(1,1)$ dan
membalikkan $(1,-1)$, sehingga pada $\mathcal B'$ matriksnya \emph{pasti}
$\operatorname{diag}(1, -1)$ --- jadi penukarannya adalah pencerminan
terhadap garis $y = x$. Adapun mencari, bagi sebuah endomorfisma yang diberikan, sebuah
\omterm{def:b1:vspaces:free}{basis} yang di situ matriksnya menjadi diagonal merupakan masalah pusat
jilid Tahun~2 (yaitu teori reduksinya); sedangkan soal akhir pekan
di bawah menunjukkan sejauh apa kesamaan \omterm{def:b1:poly:def}{polinomial} belaka sudah melangkah.
\end{example}

\begin{example}[Perubahan basis, yang dijalankan terbalik]
\label{ex:b1:matrices:reverserun}
\omterm{def:b1:linmaps:projection}{Proyeksi} pada $F = \operatorname{Vect}(1,1)$ sepanjang $G =
\operatorname{Vect}(1,-1)$ mempunyai, pada \omterm{def:b1:vspaces:free}{basis} yang disesuaikan $\mathcal
B' = \bigl((1,1),(1,-1)\bigr)$, matriks yang bening $A' =
\operatorname{diag}(1, 0)$. Untuk memperoleh matriks \omterm{def:b1:vspaces:free}{basis} kanoniknya,
jalankanlah \cref{thm:b1:matrices:conjugation} mundur, $A = P A'
P^{-1}$:
\[
P = \begin{pmatrix} 1 & 1\\ 1 & -1\end{pmatrix},
\quad
P^{-1} = \frac12\begin{pmatrix} 1 & 1\\ 1 & -1\end{pmatrix},
\quad
A = P\begin{pmatrix} 1 & 0\\ 0 & 0\end{pmatrix}P^{-1}
= \frac12\begin{pmatrix} 1 & 1\\ 1 & 1\end{pmatrix}.
\]
Periksa: $A^2 = A$ (jadi idempoten), $\operatorname{tr} A = 1 =
\operatorname{rk} A$, dan $A\binom{1}{1} = \binom11$,
$A\binom{1}{-1} = 0$, sebagaimana ditetapkan. Adapun arah terbalik ini ---
rancanglah matriksnya pada \omterm{def:b1:vspaces:free}{basis} yang baik, lalu konjugasikan kembali --- itulah
cara matriks rotasi, pencerminan dan \omterm{def:b1:linmaps:projection}{proyeksi} sesungguhnya
diproduksi dalam praktik.
\end{example}

\begin{theorem}[Bentuk normal rank]\label{thm:b1:matrices:rank}
\emph{Rank} sebuah matriks (yaitu rank kolomnya, setara dengan itu
rank \omterm{def:b1:linmaps:def}{pemetaan linear} yang berkaitan) merupakan satu-satunya invarian kesetaraannya:
karena setiap $A \in \mathcal{M}_{n,p}$ yang ber-rank $r$ setara dengan
\[
J_r = \begin{pmatrix} I_r & 0 \\ 0 & 0 \end{pmatrix},
\]
dan $\operatorname{rk}(A^{\mathsf T}) = \operatorname{rk}(A)$: jadi rank barisnya
sama dengan rank kolomnya.
\end{theorem}

\begin{proof}
Misalkan $u \colon E \to F$ ber-rank $r$. Pilihlah sebuah \omterm{def:b1:vspaces:sum}{pelengkap} $S$ bagi
$\ker u$ (dengan $\dim S = r$, \cref{thm:b1:linmaps:ranknullity}) berbasis
$(e_1, \dots, e_r)$, yang dilengkapi oleh \omterm{def:b1:vspaces:free}{basis} $\ker u$ menjadi
\omterm{def:b1:vspaces:free}{basis} $E$; maka peta $f_i = u(e_i)$, dengan $i \leq r$, membentuk \omterm{def:b1:vspaces:free}{basis}
$\operatorname{im} u$ (karena pembatasannya sebuah isomorfisma), yang dilengkapi
menjadi \omterm{def:b1:vspaces:free}{basis} $F$. Pada \omterm{def:b1:vspaces:free}{basis} inilah matriks $u$ tepat menjadi
$J_r$. Sehingga $A = Q J_r P^{-1}$ untuk $P, Q$ yang dapat dibalik.

Lalu mentransposkannya: $A^{\mathsf T} = (P^{-1})^{\mathsf T} J_r^{\mathsf T} Q^{\mathsf T}$
dengan $J_r^{\mathsf T}$ yang berbentuk sama (ber-rank $r$) dan faktor
luarnya yang dapat dibalik (karena transpos yang dapat dibalik pun dapat dibalik, dari
$(AB)^{\mathsf T} = B^{\mathsf T}A^{\mathsf T}$ yang diterapkan pada $AA^{-1} = I$):
sehingga $\operatorname{rk} A^{\mathsf T} = r$.
\end{proof}

\section{Operasi baris}

\begin{method}[Penghapusan Gauss pada matriks]\label{met:b1:matrices:gauss}
Ketiga \emph{operasi baris elementer}\index{operasi baris}nya ---
menukar dua baris, mengalikan sebuah baris dengan $\lambda \neq 0$, dan menambahkan kelipatan
sebuah baris pada baris lain --- tak mengubah ranknya (karena masing-masingnya berupa
perkalian kiri dengan matriks yang dapat dibalik). Adapun algoritmanya: ciptakanlah sebuah poros
(yaitu entri taknol yang paling kiri), bersihkanlah kolomnya di bawahnya, lalu berpindahlah ke baris
dan kolom berikutnya; dan cacah poros bentuk eselon yang dihasilkannya
adalah ranknya.

\emph{Perhitungan inversnya:} jalankanlah algoritmanya pada blok $(A \mid
I_n)$ sampai blok kirinya menjadi $I_n$ (yang mungkin jika dan hanya jika $A$
dapat dibalik); maka blok kanannya lalu menjadi $A^{-1}$ --- memang hasil kali
matriks elementer yang dipakai sama dengan $A^{-1}$.
\end{method}

\begin{example}\label{ex:b1:matrices:inverse}
$A = \begin{pmatrix} 1 & 2 \\ 3 & 4 \end{pmatrix}$: reduksikanlah $(A \mid
I_2)$:
\[
\begin{pmatrix} 1 & 2 & 1 & 0\\ 3 & 4 & 0 & 1 \end{pmatrix}
\to
\begin{pmatrix} 1 & 2 & 1 & 0\\ 0 & -2 & -3 & 1 \end{pmatrix}
\to
\begin{pmatrix} 1 & 0 & -2 & 1\\ 0 & 1 & \tfrac32 & -\tfrac12
\end{pmatrix},
\]
(dengan operasinya: $L_2 \leftarrow L_2 - 3L_1$; lalu $L_1 \leftarrow L_1
+ L_2$, $L_2 \leftarrow -\frac12 L_2$). Sehingga $A^{-1} =
\begin{pmatrix} -2 & 1 \\ \tfrac32 & -\tfrac12\end{pmatrix}$.
\emph{Periksa:} $AA^{-1} = I_2$.
\end{example}

\begin{example}[Rank dengan sebuah parameter, lewat baris belaka]
\label{ex:b1:matrices:rankparameter}
Untuk $m \in \R$, carilah rank
$M_m = \begin{pmatrix} 1 & 1 & m\\ 1 & m & 1\\ m & 1 &
1\end{pmatrix}$. Reduksikanlah: $L_2 \leftarrow L_2 - L_1$ dan $L_3
\leftarrow L_3 - mL_1$ memberikan baris
\[
(1,\ 1,\ m), \qquad (0,\ m - 1,\ 1 - m), \qquad
(0,\ 1 - m,\ 1 - m^2).
\]
\emph{Kasus $m = 1$}: kedua baris terakhirnya lenyap --- jadi satu poros,
sehingga $\operatorname{rk} M_1 = 1$ (karena ketiga baris aslinya
sama). \emph{Kasus $m \neq 1$}: skalakanlah $L_2$ dengan $\frac1{m-1}$ dan
$L_3$ dengan $\frac1{1-m}$ untuk memperoleh $(0, 1, -1)$ dan $(0, 1, 1 + m)$,
lalu $L_3 \leftarrow L_3 - L_2 = (0, 0, m + 2)$. Jika $m = -2$: maka
dua poros, jadi ber-rank $2$; kalau tidak tiga poros, jadi ber-rank $3$. Ringkasnya:
\[
\operatorname{rk} M_m =
\begin{cases}
1 & m = 1,\\
2 & m = -2,\\
3 & \text{jika tidak}.
\end{cases}
\]
Ambang yang sama akan jatuh dari satu perhitungan determinan
pada \cref{ch:b1:det} (yaitu \omterm{def:b1:poly:def}{polinomial} $-(m+2)(m-1)^2$ pada
\cref{exo:b1:det:7}) --- tetapi perhatikanlah apa yang diberikan penghapusannya yang
tak diberikan determinannya: yaitu \emph{nilai} ranknya pada
kasus yang merosot, bukan sekadar fakta bahwa ia turun.
\end{example}

\begin{example}[Menghitung pangkat]\label{ex:b1:matrices:powers}
$A = \begin{pmatrix} 1 & 1 \\ 0 & 1\end{pmatrix} = I + N$ dengan $N =
E_{12}$, $N^2 = 0$. Dan karena $I$ dan $N$ berkomutasi, teorema binomialnya
(\cref{prop:b1:structures:binomial}) terpenggal:
\[
A^k = I + kN = \begin{pmatrix} 1 & k \\ 0 & 1 \end{pmatrix}
\qquad (k \in \N, \text{ dan } k \in \Z \text{ dengan memakai } A^{-1} = I -
N).
\]
\end{example}

\begin{method}[Menghitung $A^n$: ketiga rutenya]
\label{met:b1:matrices:powers}
\begin{enumerate}
  \item \emph{Rute binomial}: jika $A = \lambda I + N$ dengan $N$ yang
        nilpoten, maka teorema binomialnya terpenggal
        (\cref{ex:b1:matrices:powers}, \cref{exo:b1:matrices:5});
        dan ia berlaku karena $\lambda I$ berkomutasi dengan segalanya.
  \item \emph{Rute \omterm{def:b1:poly:def}{polinomial}}: carilah sebuah kesamaan \omterm{def:b1:poly:def}{polinomial}
        yang dipenuhi $A$ (pada dimensi $2$, selalu $A^2 = sA -
        pI$) lalu reduksikanlah $X^n$ modulo ia; adapun soal akhir pekan
        di bawah membangun rute ini sepenuhnya.
  \item \emph{Rute keserupaan}: carilah $P$ yang dapat dibalik dengan
        $P^{-1}AP = D$ yang sederhana (yaitu diagonal, atau geseran), hitunglah
        $D^n$, lalu balikkanlah: $A^n = P D^n P^{-1}$
        (\cref{thm:b1:matrices:conjugation},
        \cref{ex:b1:matrices:conjugationrun}); adapun pencarian sistematis
        atas $P$ yang demikian adalah teori reduksi Tahun~2.
\end{enumerate}
Rute mana pun yang dipakai, periksalah hasilnya pada $n = 0, 1, 2$: yaitu tiga
uji yang murah yang menangkap hampir setiap keteledoran.
\end{method}

\begin{remark}[Jebakan yang lazim: harga ketakkomutatifannya]
\label{rem:b1:matrices:pitfalls}
Setiap kesamaan aljabar skalar yang buktinya mengurutkan ulang faktornya
mati di $\mathcal{M}_n(K)$, dengan $n \geq 2$.
\emph{Kuadratnya}: $(A + B)^2 = A^2 + AB + BA + B^2$, dan
bagian tengahnya runtuh menjadi $2AB$ hanya bila $AB = BA$
(\cref{exo:b1:matrices:1}).
\emph{Pangkat hasil kalinya}: $(AB)^k$ adalah $ABAB\cdots$, bukan
$A^kB^k$.
\emph{Pembagi nolnya}: $E_{12}E_{12} = 0$ dengan $E_{12} \neq 0$;
sehingga \emph{tak ada pencoretan}: karena $AB = AC$ mengakibatkan $B = C$
hanya ketika $A$ dapat dibalik (kalikanlah dengan $A^{-1}$ --- pada
sisi yang benar).
\emph{\omterm{def:b1:matrices:transpose}{Tracenya}}: $\operatorname{tr}(AB) = \operatorname{tr}(BA)$
selalu, tetapi $\operatorname{tr}(AB) \neq
\operatorname{tr}A\operatorname{tr}B$ pada umumnya (ambillah $A = B =
I_2$: $2 \neq 4$), dan $\operatorname{tr}(ABC) =
\operatorname{tr}(BCA)$ (secara berdaur) sedangkan $\operatorname{tr}(ACB)$
boleh berbeda.
\emph{Transposnya membalikkan}: $(AB)^{\mathsf T} = B^{\mathsf
T}A^{\mathsf T}$ --- adapun melupakan pembalikannya merupakan galat yang paling lazim
pada perhitungan keortogonalan (\cref{ch:b1:euclid}).
Jadi bila ragu, ujilah sebarang kesamaan yang diklaim pada $E_{12}$ dan
$E_{21}$: karena pasangan takkomutatif yang terkecil membantah kebanyakan rumus
yang salah dalam satu baris.
\end{remark}

\begin{remark}[Ke mana kamusnya pergi]
\label{rem:b1:matrices:whereused}
Kamus matriksnya dipakai pada setiap halaman sisa jilid
ini: \cref{ch:b1:det} melekatkan pada setiap matriks persegi satu
bilangan yang memutuskan keterbalikannya, lalu memecahkan $AX = B$
secara sistematis; \cref{ch:b1:euclid} memilih matriks
yang mengawetkan panjangnya (yaitu matriks ortogonal); sedangkan pada
\cref{ch:b1:multivar}, perilaku orde dua sebuah fungsi dua
variabel berupa matriks setangkup $2 \times 2$. Adapun \omterm{def:b1:matrices:transpose}{tracenya},
yang diperkenalkan di atas nyaris sambil lalu, menjadi invarian yang ampuh:
\cref{exo:b1:matrices:6,exo:b1:matrices:8} memberikan cicipan pertamanya,
dan jilid Tahun~2 membangun teori nilai eigen di atasnya. Sedangkan soal
akhir pekannya mengembangkan kuda beban yang lain: yaitu kesamaan \omterm{def:b1:poly:def}{polinomial}
yang dipenuhi sebuah matriks, yang mengubah perhitungan $A^n$ menjadi
\omterm{pb:b1:matrices:1}{rekurensi linear} bersuku dua.
\end{remark}

\begin{remark}[Cakrawala di dalam Buku 3]
\label{rem:b1:matrices:perspectives}
Tiga keluarga matriks yang diperkenalkan di sini mempunyai janji nanti dalam
jilid ini. \emph{Matriks setangkup}
(\cref{ex:b1:matrices:symsplit}) membawa data orde dua bagi
fungsi dua variabel: karena uji Monge pada
\cref{ch:b1:multivar} merupakan \omterm{def:b1:logic:statement}{pernyataan} tentang perilaku tanda
sebuah matriks setangkup $2\times2$, dan determinannya $rt - s^2$
dihitung oleh mesin \cref{ch:b1:det}. \emph{Matriks
ortogonal} ($A^{\mathsf T}A = I$) merupakan isometri
\cref{ch:b1:euclid}, yang di situ transposnya akhirnya memperoleh makna
geometrinya: yaitu bayangan aljabar hasil kali
dalamnya. \emph{Matriks yang dapat dibalik} bertemu uji praktisnya
pada \cref{ch:b1:det} --- yaitu satu bilangan, $\det A \neq 0$ --- yang menutup
pencarian yang dimulai bab ini dengan reduksi baris. Lalu \omterm{def:b1:matrices:transpose}{trace} dan
determinannya berjalan sebagai pasangan invarian $(s, p)$ pada
soal akhir pekannya, sampai ke teori nilai eigen
Tahun~2.
\end{remark}

\section{Latihan}

\begin{exercise}[$\star$]\label{exo:b1:matrices:1}
Misalkan $A = \begin{pmatrix} 1 & 2 \\ 0 & 1 \end{pmatrix}$ dan $B =
\begin{pmatrix} 0 & 1 \\ 1 & 0\end{pmatrix}$. Hitunglah $AB$, $BA$,
$A^2 - B^2$ dan $(A+B)(A-B)$; lalu jelaskanlah mengapa dua yang terakhir berbeda.
\end{exercise}

\begin{exercise}[$\star$]\label{exo:b1:matrices:2}
Hitunglah rank
\[
M = \begin{pmatrix}
1 & 2 & 3\\
2 & 4 & 6\\
1 & 1 & 1
\end{pmatrix},
\qquad
N = \begin{pmatrix}
1 & 1 & 0 & 2\\
0 & 1 & 1 & 1\\
1 & 2 & 1 & 3
\end{pmatrix}.
\]
\end{exercise}

\begin{exercise}[$\star$]\label{exo:b1:matrices:3}
Balikkanlah, lewat reduksi baris, $A = \begin{pmatrix} 1 & 0 & 1\\ 2 & 1 &
1\\ 1 & 1 & 1 \end{pmatrix}$, lalu periksalah pada satu hasil kali.
\end{exercise}

\begin{exercise}[$\star$]\label{exo:b1:matrices:4}
Tulislah matriks, pada \omterm{def:b1:vspaces:free}{basis} kanonik $\R_2[X]$, bagi
endomorfisma $u(P) = P(X + 1)$. Jelaskanlah, tanpa perhitungan, mengapa
ia dapat dibalik, lalu berikanlah matriks $u^{-1}$.
\end{exercise}

\begin{exercise}[$\star\star$]\label{exo:b1:matrices:5}
Misalkan $A = \begin{pmatrix} 2 & 1 \\ 0 & 2\end{pmatrix}$. Tulislah $A =
2I + N$, hitunglah $N^2$, lalu simpulkanlah $A^k$ untuk setiap $k \in \N$ lewat
teorema binomialnya.
\end{exercise}

\begin{exercise}[$\star\star$]\label{exo:b1:matrices:6}
Buktikan bahwa tak ada matriks $A, B \in \mathcal{M}_n(K)$ (dengan
$K = \R$ atau $\C$) sedemikian sehingga $AB - BA = I_n$. \emph{(Ambillah \omterm{def:b1:matrices:transpose}{tracenya}.)}
\end{exercise}

\begin{exercise}[$\star\star$]\label{exo:b1:matrices:7}
Sebuah matriks $A$ disebut \emph{nilpoten} bila $A^m = 0$ untuk suatu $m$. Buktikan
bahwa $I - A$ lalu dapat dibalik, dengan
\[
(I - A)^{-1} = I + A + A^2 + \dots + A^{m-1} .
\]
Penerapannya: balikkanlah $\begin{pmatrix} 1 & 2 & 3\\ 0 & 1 & 2\\ 0 & 0 &
1\end{pmatrix}$.
\end{exercise}

\begin{exercise}[$\star\star$]\label{exo:b1:matrices:8}
Misalkan $A \in \mathcal{M}_n(\R)$ memenuhi $A^2 = A$ (yaitu idempoten). Buktikan
bahwa $\operatorname{tr} A = \operatorname{rk} A$. \emph{(Tafsirkanlah
$A$ sebagai sebuah \omterm{def:b1:linmaps:projection}{proyeksi} lalu pilihlah \omterm{def:b1:vspaces:free}{basis} yang disesuaikan;
adapun \cref{thm:b1:matrices:conjugation} mengatakan bahwa \omterm{def:b1:matrices:transpose}{tracenya}
tak bergantung pada basisnya karena $\operatorname{tr}(P^{-1}MP) =
\operatorname{tr} M$.)}
\end{exercise}

\begin{exercise}[$\star\star\star$]\label{exo:b1:matrices:9}
Misalkan $J \in \mathcal{M}_n(\R)$ matriks serba satu. Hitunglah $J^2$,
lalu simpulkanlah, untuk $a, b \in \R$, syarat keterbalikan $M
= aI + bJ$ beserta $M^{-1}$ \emph{(carilah invers yang
berbentuk sama $\alpha I + \beta J$)}.
\end{exercise}

\begin{exercise}[$\star\star\star$]\label{exo:b1:matrices:10}
(Ketaksamaan rank) Untuk $A, B \in \mathcal{M}_n(K)$, buktikanlah
\[
\operatorname{rk}(A + B) \leq \operatorname{rk} A +
\operatorname{rk} B,
\qquad
\operatorname{rk}(AB) \geq \operatorname{rk} A + \operatorname{rk}
B - n .
\]
\emph{(Untuk yang kedua --- yaitu ketaksamaan Sylvester --- terapkanlah
rank--nulitas pada pembatasan \omterm{def:b1:logic:map}{pemetaan} $A$ ke
$\operatorname{im} B$.)}
\end{exercise}

\begin{exercise}[$\star\star$]\label{exo:b1:matrices:11}
Misalkan $D = \operatorname{diag}(d_1, \dots, d_n)$ dengan $d_i$ yang
\emph{berbeda berpasangan}.
\begin{enumerate}
  \item Buktikan bahwa sebuah matriks $A$ berkomutasi dengan $D$ jika dan hanya jika
        $A$ diagonal. \emph{(Bandingkanlah entri $(i,j)$
        $AD$ dan $DA$.)}
  \item Simpulkanlah \emph{pusat} $\mathcal{M}_n(K)$: bahwa
        matriks yang berkomutasi dengan \emph{setiap} matriks tepat berupa
        matriks skalar $\lambda I_n$. \emph{(Ujilah terhadap
        $D$, lalu terhadap matriks $E_{ij}$.)}
\end{enumerate}
\end{exercise}

\begin{exercise}[$\star\star\star$]\label{exo:b1:matrices:12}
(Matriks ber-rank satu) Misalkan $A \in \mathcal{M}_n(K)$, $A \neq 0$.
\begin{enumerate}
  \item Buktikan bahwa $\operatorname{rk} A = 1$ jika dan hanya jika $A =
        CL$ untuk suatu kolom taknol $C \in \mathcal{M}_{n,1}$ dan
        baris taknol $L \in \mathcal{M}_{1,n}$.
  \item Untuk $A$ yang demikian, buktikanlah $A^2 = (\operatorname{tr} A)\,A$;
        lalu simpulkanlah bahwa matriks ber-rank satu bersifat nilpoten jika dan hanya jika
        \omterm{def:b1:matrices:transpose}{tracenya} nol.
  \item Jika $\operatorname{tr} A \neq -1$, buktikanlah bahwa $I_n + A$
        dapat dibalik dengan
        \[
        (I_n + A)^{-1} = I_n - \frac{1}{1 + \operatorname{tr}
        A}\,A ,
        \]
        dan bahwa $I_n + A$ \emph{tak} dapat dibalik ketika
        $\operatorname{tr} A = -1$. \emph{(Carilah sebuah vektor yang dibunuh
        oleh $I_n + A$.)}
\end{enumerate}
\end{exercise}

\section{Soal: pangkat sebuah matriks lewat pembagian polinomial}

\begin{problem}\label{pb:b1:matrices:1}
Menghitung $A^{100}$ entri demi entri tanpa harapan; sedangkan menghitungnya
lewat sebuah kesamaan \omterm{def:b1:poly:def}{polinomial} yang dipenuhi $A$ memakan tiga baris.
Soal ini membangun metodenya dari nol: yaitu pembagian Euclid
atas $X^n$, kesamaan $A^2 - sA + pI = 0$ yang dipenuhi setiap matriks $2
\times 2$ (yakni teorema Cayley--Hamilton dalam dimensi
$2$), dan kamus antara pangkat matriks dan \omterm{pb:b1:matrices:1}{rekurensi
linear} --- dengan \omterm{pb:b1:matrices:1}{bilangan Fibonacci} sebagai contoh yang berjalan.
\index{teorema Cayley--Hamilton (dimensi 2)}
\index{bilangan Fibonacci}\index{rekurensi linear}

\medskip
\noindent\textbf{Bagian I --- Kalkulus sisanya.} Tetapkanlah $s, p
\in K$ dan $D = X^2 - sX + p$.
\begin{enumerate}
  \item Benarkanlah bahwa untuk setiap $n \in \N$ ada $Q_n
        \in K[X]$ dan $(a_n, b_n) \in K^2$ yang tunggal dengan
        \[
        X^n = Q_n\,D + a_n X + b_n ,
        \]
        lalu hitunglah $(a_0, b_0)$ dan $(a_1, b_1)$.
  \item Dengan mengalikannya dengan $X$ lalu \omterm{def:b1:arith:divides}{membaginya} lagi, tegakkanlah
        rekurensinya
        \[
        a_{n+1} = s\,a_n + b_n,
        \qquad
        b_{n+1} = -p\,a_n ,
        \]
        lalu simpulkanlah $a_{n+2} = s\,a_{n+1} - p\,a_n$: sehingga
        barisan koefisiennya menuruti \omterm{pb:b1:matrices:1}{rekurensi linear} yang melekat
        pada $D$.
  \item Andaikan $D$ berakar dua yang berbeda $\lambda \neq \mu$.
        Dengan menilai kesamaan pembagiannya, buktikanlah
        \[
        a_n = \frac{\lambda^n - \mu^n}{\lambda - \mu},
        \qquad
        b_n = \frac{\lambda\mu^n - \mu\lambda^n}{\lambda - \mu} .
        \]
  \item Andaikan $D = (X - \lambda)^2$. Dengan memakai \omterm{def:b1:derivative:def}{turunan}
        kesamaan pembagiannya, buktikanlah $a_n = n\lambda^{n-1}$ dan
        $b_n = (1 - n)\lambda^{n}$.
  \item Tunjukkan bahwa menyulihkan sebuah matriks tetap $M \in
        \mathcal{M}_k(K)$ ke dalam \omterm{def:b1:poly:def}{polinomial} menghormati jumlah dan
        hasil kalinya: $(PQ)(M) = P(M)\,Q(M)$. Lalu simpulkanlah bahwa jika $D(M) =
        0$, maka
        \[
        M^n = a_n\,M + b_n\,I \qquad (n \in \N).
        \]
\end{enumerate}

\medskip
\noindent\textbf{Bagian II --- Dimensi 2: \omterm{def:b1:matrices:transpose}{trace}, bilangan
determinan, Cayley--Hamilton.} Untuk $A = \begin{pmatrix} a & b\\ c &
d\end{pmatrix}$ tetapkanlah $s = a + d = \operatorname{tr} A$ dan $p = ad
- bc$ (yaitu bilangan yang akan dinamai \cref{ch:b1:det} sebagai
determinan).
\begin{enumerate}[resume]
  \item Periksalah lewat perhitungan langsung \emph{kesamaan
        Cayley--Hamilton dalam dimensi $2$}:
        \[
        A^2 - s\,A + p\,I_2 = 0 .
        \]
  \item Buktikan lewat penjabaran langsung bahwa $p$ bersifat perkalian:
        yakni, dengan notasi yang kasatmata, $p(AB) = p(A)\,p(B)$. Lalu tunjukkan:
        bahwa $A$ dapat dibalik jika dan hanya jika $p \neq 0$, dan pada
        kasus itu
        \[
        A^{-1} = \frac1p\,\bigl(s\,I_2 - A\bigr).
        \]
  \item Misalkan $A = \begin{pmatrix} 1 & 1\\ 0 & 2\end{pmatrix}$.
        Hitunglah $s$, $p$, akar $D$, lalu simpulkanlah rumus
        \omterm{def:b1:topology:closed}{tertutup} bagi $A^n$; lalu periksalah terhadap perhitungan langsung
        atas $A^2$.
  \item Misalkan $A = \begin{pmatrix} 3 & 1\\ -1 & 1\end{pmatrix}$.
        Tunjukkan bahwa $D$ berakar rangkap lalu hitunglah $A^n$; lalu periksalah
        di $n = 2$.
  \item Misalkan $F = \begin{pmatrix} 1 & 1\\ 1 & 0\end{pmatrix}$ lalu
        definisikanlah \omterm{pb:b1:matrices:1}{bilangan Fibonacci} lewat $F_0 = 0$, $F_1 = 1$,
        $F_{n+2} = F_{n+1} + F_n$. Buktikanlah
        \[
        F^n = \begin{pmatrix} F_{n+1} & F_n\\ F_n &
        F_{n-1}\end{pmatrix} \quad (n \geq 1),
        \]
        lalu simpulkanlah rumus Binet $F_n = \dfrac{\varphi^n -
        \psi^n}{\sqrt5}$ dengan $\varphi = \frac{1 + \sqrt5}2$,
        $\psi = \frac{1 - \sqrt5}2$, dan, dengan memakai pertanyaan 7,
        kesamaan Cassini $F_{n+1}F_{n-1} - F_n^2 = (-1)^n$.
\end{enumerate}

\medskip
\noindent\textbf{Bagian III --- \omterm{pb:b1:matrices:1}{Rekurensi linear}, secara struktural.}
Tetapkanlah $s, p \in K$ dengan $p \neq 0$, lalu misalkan $E_D$ \omterm{def:b1:logic:sets}{himpunan}
barisan dengan $u_{n+2} = s\,u_{n+1} - p\,u_n$ untuk setiap $n$.
\begin{enumerate}[resume]
  \item Tunjukkan bahwa $E_D$ merupakan \omterm{def:b1:vspaces:def}{ruang vektor} berdimensi $2$
        (sesuaikanlah \cref{exo:b1:findim:10}).
  \item Tunjukkan bahwa barisan $(a_n)$ pada Bagian I merupakan unsur
        $E_D$ dengan nilai awalnya $0, 1$, dan bahwa setiap $u
        \in E_D$ memenuhi
        \[
        u_n = u_1\,a_n + u_0\,b_n \qquad (n \in \N),
        \]
        dengan $(b_n)$ seperti pada Bagian I: sehingga sisa pembagiannya memecahkan
        \emph{semua} rekurensinya sekaligus.
  \item Jika $\lambda \neq \mu$ akar $D$, tunjukkanlah bahwa
        $\bigl((\lambda^n), (\mu^n)\bigr)$ merupakan \omterm{def:b1:vspaces:free}{basis} $E_D$;
        sedangkan jika $D = (X-\lambda)^2$ dengan $\lambda \neq 0$, tunjukkanlah bahwa
        $\bigl((\lambda^n), (n\lambda^n)\bigr)$ merupakan satu.
  \item Pecahkanlah selengkapnya: $u_{n+2} = u_{n+1} + 6u_n$, $u_0 = 1$,
        $u_1 = 8$; lalu periksalah jawabannya pada $u_2$ dan $u_3$.
  \item Misalkan $C = \begin{pmatrix} 0 & 1\\ -p & s\end{pmatrix}$
        (yaitu \emph{matriks pendamping} $D$). Tunjukkan bahwa
        \[
        \begin{pmatrix} u_{n}\\ u_{n+1}\end{pmatrix}
        = C^n \begin{pmatrix} u_0\\ u_1\end{pmatrix}
        \quad (u \in E_D),
        \]
        dan bahwa $\operatorname{tr} C = s$ serta $p(C) = p$: sehingga
        rekurensi dan matriksnya membawa \omterm{def:b1:poly:def}{polinomial} $D$ yang sama.
\end{enumerate}

\medskip
\noindent\textbf{Bagian IV --- Derajat tiga.} Misalkan $D_3 = X^3 -
\alpha X^2 - \beta X - \gamma$ dan
\[
C_3 = \begin{pmatrix} 0 & 1 & 0\\ 0 & 0 & 1\\ \gamma & \beta &
\alpha \end{pmatrix}.
\]
\begin{enumerate}[resume]
  \item Tunjukkan bahwa $D_3(C_3) = 0$. \emph{(Hitunglah peta
        vektor \omterm{def:b1:vspaces:free}{basis} kanoniknya di bawah pangkat $C_3$: karena
        \omterm{def:b1:logic:map}{pemetaan} $C_3$ mengirim $e_1 \mapsto \dots \mapsto$ sebuah
        kombinasi yang dipaksa oleh baris terakhirnya.)}
  \item Tunjukkan bahwa jika $D_3$ berakar tiga yang berbeda $\lambda_1,
        \lambda_2, \lambda_3$, maka sisa $R_n$ atas $X^n$
        yang dibagi $D_3$ merupakan \emph{interpolan Lagrange} atas
        nilai $\lambda_i^n$ pada simpul $\lambda_i$
        (\cref{thm:b1:poly:lagrange}); lalu simpulkanlah bahwa setiap entri
        $C_3^{\,n}$ merupakan kombinasi \omterm{def:b1:linmaps:def}{linear} yang tetap atas
        $\lambda_1^n, \lambda_2^n, \lambda_3^n$.
  \item Pecahkanlah: $u_{n+3} = 2u_{n+2} + u_{n+1} - 2u_n$ dengan $u_0 =
        0$, $u_1 = 1$, $u_2 = 1$. \emph{(Faktorkanlah $D_3 = (X - 1)(X
        + 1)(X - 2)$.)} Lalu periksalah pada $u_3$.
  \item Hitunglah sisa $X^n$ modulo $(X -
        \lambda)^3$ \emph{(lewat ekspansi Taylor $X^n$ di
        $\lambda$)}, lalu simpulkanlah sebuah rumus bagi $(\lambda I +
        N)^n$ ketika $N^3 = 0$ dan $N$ berkomutasi dengan segalanya yang
        terlibat; lalu periksalah ia terhadap teorema binomialnya.
  \item Tunjukkan bahwa untuk $D_3$ yang berakar berbeda, penyelesaian umum
        rekurensi berorde $3$-nya adalah $u_n = c_1
        \lambda_1^n + c_2\lambda_2^n + c_3\lambda_3^n$: yaitu buktikanlah
        bahwa ketiga barisan geometrinya membentuk \omterm{def:b1:vspaces:free}{basis}
        ruang penyelesaiannya. \emph{(Untuk kebebasannya, nilailah sebuah
        kombinasi nol di $n = 0, 1, 2$ lalu kenalilah sebuah
        sistem interpolasi pada simpul yang berbeda
        $\lambda_i$.)}
\end{enumerate}

\medskip
\noindent\textbf{Bagian V --- Dividen Fibonacci, dan sintesisnya.}
\begin{enumerate}[resume]
  \item Buktikan $F_1 + F_2 + \dots + F_n = F_{n+2} - 1$.
  \item Dari $F^{m+n} = F^m F^n$, turunkanlah rumus penjumlahannya
        \[
        F_{m+n} = F_{m+1}F_n + F_m F_{n-1},
        \]
        lalu simpulkanlah $F_{2n} = F_n(F_{n+1} + F_{n-1})$.
  \item Buktikan bahwa $F_n$ merupakan bilangan bulat terdekat dengan
        $\varphi^n/\sqrt5$ untuk setiap $n \geq 0$.
  \item Misalkan $t_n = \operatorname{tr}(F^n) = F_{n+1} + F_{n-1}$
        (yaitu \emph{bilangan Lucas} $L_n$). Tunjukkan $t_{n+2} =
        t_{n+1} + t_n$, $t_1 = 1$, $t_2 = 3$, bahwa $L_n =
        \varphi^n + \psi^n$, lalu pulihkanlah $F_{2n} = F_n L_n$.
  \item Sintesis, dalam empat kalimat: mengapakah pangkat matriks $2
        \times 2$ hidup di bidang
        $\operatorname{Vect}(I, A)$ pada $\mathcal{M}_2(K)$ (yaitu argumen
        dimensi yang mana yang menjamin sebuah kesamaan kuadrat, dan
        kesamaan eksplisit yang mana yang dihasilkan Bagian II); bagaimanakah pembagian
        Euclid mengubah pemangkatan menjadi rekurensi
        bersuku dua; \omterm{def:b1:logic:statement}{pernyataan} yang mana pada soal ini yang merupakan kasus $n =
        2$ sebuah teorema yang sah pada segala dimensi (namailah ia,
        lalu katakanlah di mana ia dibuktikan pada seri ini); dan apa yang
        ditambahkan konstruksi matriks pendampingnya pada gambarannya.

\end{enumerate}
\end{problem}
